\subsection*{6.2.1 Proof of Theorem 5}
In this short section we prove Theorem 5. We focus on the version that is based on non-explicit samplers, yielding non-explicit reductions. The explicit reduction, which entails a bit more technical work, is deferred to Section 6.3 and Section 6.6. We choose to focus on the non-explicit version first because we believe that understanding LDC in the information-theoretic level is, at present, a deeper and more urgent problem than the question of explicitness. Also, the parameters are easier to work with. For the information-theoretic version, we make use of the sampler that is given by Theorem 15. From here on we refer to the constant $c_{\mathrm{samp}}\geq1$ that appears in that theorem.
\begin{proof}[Author proof of Theorem 72, stated at the start of this item]
Take $C_1$ to be the code $C$ in the hypothesis of the theorem, namely, a code with block length $n_1=n$ and distance $\delta_1=\delta$. Recall that in the distance amplification procedure from Section 6.1, we make use of a $(\delta_2/2,\delta_1)$ sampler $G=([\ell],[r],E)$ with $\ell=n_1/k_2$ and left-degree $d=n_2$. For the proof, we will instantiate the distance amplification procedure with the sampler that is given by Theorem 15. We take $C_2$ to be an asymptotically good code over $\Sigma$ set with block length
\[n_2=c_{\mathrm{samp}}\cdot\frac{\log(1/\delta_1)}{(\delta_2/2)^2}=O(\log(1/\delta)),\]
where $\delta_2$ is the (constant) distance of $C_2$, having rate at least $1/2$. Note that this choice of parameters is as required by Theorem 15 from the left degree of the sampler. Clearly, $C_2$ has query complexity $O(\log(1/\delta))$ and error $\varepsilon_2=0$. As for the degree $D_v$ of any given right vertex $v$ of the sampler, note that the average right degree is
\[D=\frac{\ell d}{r}=\frac d{\delta\log(1/\delta)}=\frac{4c_{\mathrm{samp}}}{\delta_1\delta_2^2}=\Theta\left(\frac1\delta\right).\]
Recall that, by Theorem 15, $D_v\in[D/2,2D]$. For every length in this range, we take a code from the family $C'$ having the required message length. We would like take the family of codes $C_3$ to be $C'$ though we must reduce the error first. Indeed, note that the error $\varepsilon$ of the code obtained by Proposition 65 is $\varepsilon_1+n(\varepsilon_2+\varepsilon_3)$. As mentioned in the introduction, one can reduce the error from $1/5$ to $1/(10n)$ by applying the decoding procedure for $c\log n$ times, where $c$ is some large enough constant, and output the symbol according to plurality. This has no effect on the rate or distance of $C'$, and has a multiplicative $O(\log n)$ cost in query complexity. That is, the query complexity of $C_3$ is $O(q'_{O(1/\delta_1)}\log n)$. The proof then readily follows by Proposition 65.
\end{proof}
\subsubsection*{6.2.1.1 Improving the query complexity further given low-error LDC}
We remark that, if $C'$ has error $O(1/n)$ to begin with, $n$ being the block length of $C$, then one can skip the error reduction in the proof of Theorem 72, and get a slightly better query complexity. Indeed, this will save the $\log n$ factor in Equation (6.2). Moreover, observe that $C_2$ can be taken to be an LDC as well, rather than a standard code, which will reduce its deterioration on the query complexity from $O(\log(1/\delta))$ to $q'_{O(\log(1/\delta))}$. However, for that, one need the error of $C_2$ to be $O(1/n)$ as well. Assuming one can obtain such low-error LDC (note that an error of $1/n$ is at least exponentially-small in the length of $C_2$ since $\delta>1/n$), the query complexity can be improved further to
\[q_{\mathrm{new}}\leq q\cdot q'_{O(1/\delta)}q'_{O(\log(1/\delta))}.\]
We conclude this section by instantiating Theorem 72 with $C'$ taken to be the state-of-the-art construction of asymptotically good LDC.
\setcounter{nttheorem}{72}
\begin{nttheorem}[Author Theorem 73, {\cite{ref26}}]
Let $\Sigma$ be a finite alphabet. Then, there exist constants $\delta,\rho$ and an explicit infinite family of $(q_k,\delta,1/5)$-LDC, $k$ being the message length, having query complexity $q_k=2^{O(\sqrt{\log(k)\log\log k})}$.
\end{nttheorem}
Using it, one gets query complexity
\[q_{\mathrm{new}}\leq q\log(n)\cdot2^{O(\sqrt{\log(1/\delta)\cdot\log\log(1/\delta)})}=q\log(n)(1/\delta)^{o(1)}.\]
